% Independent contribution. AI-produced by GPT-6 Astra; no priority claimed.
\subsection{Cyclic sphere quotients and the physical modular descent}

We use the full centralizer topology proved in Part A.1, the actual
ordinary genus-zero GT action, and (A.K2).  The construction below is
on punctured spheres throughout.

\begin{result}{B.1.7}{Proved}{The parameter on a cyclic prime-power sphere quotient}
Let $N=p^a$ and let $C_N$ rotate a sphere with its two fixed points
marked and with three individually named regular marked orbits.
The action induced by any $F\in\widehat{GT}$ on the completed
centralizer quotient
\[
 \widehat C/C_N\simeq\widehat\Gamma_{0,5}
\]
has the same full GT parameter $F$.  This statement uses the two
fixed points and the three orbit names as the five quotient labels.
The comparison is compatible with forgetting any regular orbit and
with the choice subgroups used below.
\end{result}
\begin{proof}
Write the fixed points as $0,\infty$, and write the regular orbits
as $O_i=z_i\mu_N$, $i=0,1,2$.  The generator $h$ sends
$z_i\zeta_N^j$ to $z_i\zeta_N^{j+1}$.  Let $C$ be the subgroup of
the discrete centralizer which fixes the two fixed labels and each
orbit as a set.  The cyclic cover has monodromy
\[
 d_0\mapsto h,\qquad d_\infty\mapsto h^{-1},\qquad
 d_i\mapsto1\quad(i=0,1,2).                             \tag{B.K1}
\]
Every pure mapping class of the five-punctured quotient preserves
this homomorphism, and hence lifts to a mapping class commuting
with $h$.  The kernel consists exactly of the deck transformations.
Thus the discrete covering map gives
\[
             1\longrightarrow C_N\longrightarrow C
                \longrightarrow\Gamma_{0,5}\longrightarrow1.
                                                               \tag{B.K2}
\]
All branch points are marked.  The centralizer full-topology proof
of Part A.1 therefore applies directly.  In particular the actual
closed centralizer is $\widehat C$, and quotienting by the finite
subgroup gives the full group $\widehat\Gamma_{0,5}$ in (B.K2).

Here are the label and generator corrections, including their
multiplier.  Put $F=(\lambda,f)$.  Prime-power finite-subgroup
conjugacy first makes the image of $C_N$ discrete.  As in the
pure-label correction of Part A.1, its initial conjugator can be
made pure while allowing the resulting discrete action to vary.
The resulting cyclic action has the same labelled orbit
permutations, with fixed labels $0,\infty$.  The inverse image of
$C_N$ in the ordinary universal-curve extension is the actual
fundamental group of the quotient punctured sphere:
\[
 1\longrightarrow\widehat F_{3N+1}
  \longrightarrow\widehat F_4\longrightarrow C_N\longrightarrow1.
                                                               \tag{B.K3}
\]
This is a finite-index comparison, so no procongruence completion
is substituted.  A peripheral at an upstairs fixed puncture maps
to $d_0^N$.  The universal-curve GT action carries this upstairs
peripheral to its $\lambda$th power, up to conjugacy inside the
surface kernel.  Consequently the induced isomorphism of the
middle free groups carries $d_0^N$ to a conjugate of
$(d'_0)^{N\lambda}$.  The power-centralizer assertion used here has an elementary
cover proof.  If $g$ centralizes $d_0^N$ but does not belong to
$I=\langle d_0\rangle$, choose an open subgroup $U$ containing
$I$ and not $g$.  Its finite punctured-sphere cover has at least
three cusps.  The cusp of local degree one determined by $I$
and the distinct cusp determined by the coset $Ug$ give two
independent peripheral lines in $U^{\rm ab}$.  If the latter
cusp has local degree $e$, the equation
$g d_0^N g^{-1}=d_0^N$ forces $e\mid N$ and identifies the
nonzero multiples $N[d_0]$ and $(N/e)[d'_0]$ of those distinct
lines, a contradiction.  A discrete representative of $Ug$
identifies the latter line with an ordinary cusp; the remaining
factor of $g$ lies in $\widehat U$ and has no effect on
abelianization.  Thus $C(d_0^N)=I$.  Applying this to the
preceding conjugate power shows that the image of $d_0$ is a
conjugate of $(d'_0)^\lambda$: multiplication by $N$ is injective
on $\widehat{\mathbb Z}$.  Thus the labelled monodromy at $0$ of
the new action is $h^{\lambda^{-1}}$.

The classification here can be read directly from the cyclic
cover.  Its monodromy at the two fixed punctures is $(h^b,h^{-b})$,
with $b$ a unit modulo $N$, and at every other puncture it is one.
These data determine the oriented covering up to a homeomorphism
fixing the quotient labels.  The preceding computation therefore
supplies a geometric correction taking the new generator to
$h^\lambda$.  On each labelled regular orbit its permutation is
\[
             z_i\zeta_N^j\longmapsto z_i\zeta_N^{\lambda j};
                                                               \tag{B.K4}
\]
it fixes $0,\infty$ and every representative $z_i$ with $j=0$.
To fix all these representatives simultaneously, first fix one by a
deck translate.  Moving any other regular quotient point once around
$0$ changes its chosen sheet by $h$ and leaves the other sheets
unchanged.  These pure quotient motions adjust the remaining
representatives independently, giving exactly (B.K4).
After this one correction, the action preserves $\widehat C$, its
quotient, and all the choice subgroups fixing these representatives.
For $\lambda=1$ this correction is pure.  Applying ordinary
forgetting to this single correction provides compatible
corrections after any entire orbit is forgotten.

By the unrestricted outer automorphism theorem for
$\widehat\Gamma_{0,5}$
\cite[Corollary 2.8]{CA_HMM}, its induced outer automorphism has
form $(\Theta,\sigma)\in\widehat{GT}\times\mathfrak S_5$.
The three kernels which forget one regular orbit are preserved,
because they come from the ordinary upstairs orbit-forgetting
maps.  Hence $\sigma$ fixes all three orbit labels and can only
interchange $0,\infty$.

We now identify both $\sigma$ and $\Theta$ by actual maps.  Forget
$O_2$.  Stabilizing the representative $z_0$ gives a subgroup
$C^0$ mapping isomorphically to $C/C_N$: each deck orbit has
exactly one lift fixing $z_0$.  Stabilizing $z_1$ as well gives an
index-$N$ subgroup $C^{00}$ of $C^0$.  Normalize $z_0=1$, put
$z=z_1/z_0$ and $u=z^N$.  The three maps are precisely
\[
 \begin{array}{ccc}
 \mathbb P^1_z\setminus\{0,\mu_N,\infty\}
       &\xrightarrow{\ z\mapsto z\ }&
           \mathbb P^1_z\setminus\{0,1,\infty\}\\
 \big\downarrow {z\mapsto z^N}&&\\[-2mm]
 \mathbb P^1_u\setminus\{0,1,\infty\}.&&
 \end{array}                                                   \tag{B.K5}
\]
On full fundamental groups the vertical map is $H_N\subset
\widehat F(x,y)$, and the horizontal map is literally $\delta_N$
of (A.K1).  Upstairs the horizontal map is the ordinary forgetting
which retains exactly $0,\infty,z_0,z_1$.  The correction (B.K4)
fixes these four labels, so the action on its target has parameter
$F$, with only the usual inner basepoint discrepancy.

If $\sigma$ interchanged the two fixed labels, the induced action
on the cover would interchange its peripherals at $z=0,\infty$.
Both survive the horizontal map in (B.K5).  Their images have
nonproportional primitive vectors $(1,0)$ and $(-1,-1)$ in the
target abelianization.  This contradicts the named pure GT action
on that target, which preserves each peripheral with its unit
multiplier.  Thus $\sigma=1$.

The normalized action induced through the horizontal map is now
exactly $D_N(\Theta)$: if $\Theta=(\mu,g)$, its conjugating word
is $\delta_N(g)$, whose abelianization is $(0,c_N(g))$;
multiplying on the left by $y^{-c_N(g)}$ leaves the automorphism
unchanged and normalizes this word.  Ordinary forgetting in
(B.K5) therefore says
\[
                         D_N(\Theta)=F.                \tag{B.K6}
\]
There is no lost outer ambiguity, since the normalized gentle
embedding into $\operatorname{Out}(\widehat F_2)$ is faithful.
But $\Theta$ is an actual full GT element, so (A.K2) gives
$D_N(\Theta)=\Theta$.  Hence $\Theta=F$, proving the assertion
and its stated compatibility.
The same assertion holds with more than three regular orbits:
HMM applies to the quotient $\widehat\Gamma_{0,m+2}$, and
forgetting to any three named regular orbits restricts its GT
factor to the factor just identified.  That restriction is
faithful on $\widehat{GT}$.  The individual orbit-forgetting
kernels and the same two retained branch cusps pin the labels.
Configurations with fewer regular orbits are obtained by adding
regular orbits, making this compatible correction once, and then
using the actual ordinary forgetting map.
\end{proof}

\paragraph{The equianharmonic opening, including its marking.}
Take $N=3$ in the preceding construction and forget $O_2$.
Use instead the choice subgroup which fixes the representative
of $O_1$ and allows the three labels of $O_0$ to rotate.
It is again isomorphic to $C/C_3=\Gamma_{0,4}$.  Forget the two
fixed branch labels and the two unchosen labels of $O_1$.
The target retains the three labels of $O_0$ and the chosen
moving point.  Its mixed four-point group is the actual
index-two subgroup
\[
 R_E=\langle x,r\rangle\subset B=\widehat{C_2*C_3},
 \qquad r^3=(xr)^3=1,\qquad x=\tau^{-2}.                \tag{B.K7}
\]
One can check the orientation and branch assignment as follows.
On the four-point coordinate line $t$, the order-three map
$t\mapsto(1-t)^{-1}$ permutes $0,1,\infty$ and fixes
$\varrho=e^{\pi i/3},\varrho^{-1}$.  The coordinate
\[
 w=\frac{t-\varrho}{t-\varrho^{-1}},\qquad u=w^3          \tag{B.K8}
\]
sends these fixed points to $0,\infty$ and the three ordinary
punctures to $\mu_3$.  The quotient has ordinary cusp $u=1$
and order-three cone points $u=0,\infty$.  Choose its positive
meridian at $u=1$ to be $x$; the positive meridian at $u=0$ is
$r$, which cycles the three upstairs peripheral classes
$x,y,z$.  The remaining meridian is $(xr)^{-1}$.  Equivalently,
in the positive braid presentation $B_3=\langle a,b\mid aba=bab\rangle$,
put $r=ab$, $s=aba$, and quotient by $(ab)^3$; then
$a=\tau^{-1}$, $x=a^2$, and these are precisely (B.K7).
Changing the ordering of the three source cusps is the displayed
geometric $\mathfrak S_3$ relabelling, which commutes with the
pure GT factor in the outer action.

Thus the ordinary forget map, after passing through the choice
subgroup, is the actual surjection
\[
 q_E:\widehat F(X,Y)\twoheadrightarrow R_E,\qquad
              X\mapsto x,\quad Y\mapsto r.             \tag{B.K9}
\]
Its kernel is the closed normal subgroup generated by
$Y^3,(XY)^3$.  For multiplier one the correction (B.K4) fixes every
upstairs label.  The cyclic comparison just proved identifies
its source parameter with $F$, while ordinary forgetting
identifies its target with the restriction of the genuine
modular action of $F$.  In particular its square is
outer-commutative for the same multiplier-one $F$ on both sides.

\begin{result}{B.1.8}{Proved}{Physical descent fixes full profinite GT}
For the physical quotient $\varphi_L(x)=\tau^{-1}$,
$\varphi_L(y)=s$ of (B.C1)--(B.C3),
\[
             D_L|_{\widehat{GT}}=\operatorname{id},
 \qquad D_P(F)=cFc\quad(F\in\widehat{GT}).               \tag{B.K10}
\]
Consequently $\widehat{GT}\subseteq C_L$.  Every character
$\Delta\circ D_P^n$, $n\geq1$, vanishes on full GT.
\end{result}
\begin{proof}
First take $F=(1,f)$.  Write its genuine normalized modular
extension as
\[
 \alpha(\tau)=\tau,\qquad \alpha(r)=A rA^{-1},\qquad A\in\widehat F(x,y).
\]
The square (B.K9) has one inner discrepancy.  Since both
normalized source and target fix $x$, this discrepancy is
$x^t$ for some $t\in\widehat{\mathbb Z}$.  Indeed a nontrivial
residue in $R_E/\widehat F(x,y)=C_3$ cycles the three distinct
peripheral homology lines of the pure free group, so an element
centralizing $x$ has trivial residue.  Within that pure group,
the primitive peripheral centralizer is $\langle x\rangle$,
by the preceding finite-cover argument with $N=1$.
Comparison at $Y$ gives
\[
 f(x,r)^{-1}rf(x,r)=x^t A rA^{-1}x^{-t}.
\]
The word $f(x,r)x^tA$ is pure and centralizes $r$.
But $C_{R_E}(r)=\langle r\rangle$, whose intersection
with the pure subgroup is trivial.  Therefore
\[
                         A^{-1}=f(x,r)x^t.              \tag{B.K11}
\]
This is the equianharmonic primitive comparison with the
actual frame $A$.  Its abelian exponent need not be named.

Put $a=\tau^{-1}$, so $x=a^2$.  The already proved exact
relation (A.J20) says
$f(a,r)=r^{-\rho}f(x,r)a^{2\rho}$.
Substitution in (B.K11) gives
\[
 A^{-1}=r^\rho f(a,r)a^{2t-2\rho},\qquad
 \alpha(r)=a^{-(2t-2\rho)}f(a,r)^{-1}rf(a,r)a^{2t-2\rho}.
                                                               \tag{B.K12}
\]
Let $\psi$ be the actual automorphism descended through
$\varphi_L$ before the final inner normalization.  It fixes
$a$ and sends $s$ to $f(a,s)^{-1}s f(a,s)$.  The exact
source change $\theta'(X)=X$, $\theta'(Y)=(XY)^{-1}$ of
(A.J8) commutes with the genuine normalized multiplier-one
action.  Consequently it sends the third peripheral
$(XY)^{-1}$ to its conjugate by
$f(X,(XY)^{-1})$.  Applying the actual quotient
$X\mapsto a,Y\mapsto s$, for which $(as)^{-1}=r$, gives
\[
                         \psi(r)=f(a,r)^{-1}rf(a,r).
\]
Equation (B.K12) now identifies the two automorphisms on the
actual generating pair $(a,r)$ of $B$:
\[
                     \alpha=\operatorname{Inn}(a^{-(2t-2\rho)})\psi.
                                                               \tag{B.K13}
\]
The inner element in (B.K13) is the pure element
$x^{\rho-t}$.  Their restrictions to the pure free subgroup
therefore have the same outer class.  Faithfulness of the normalized gentle action
therefore gives $D_L(F)=F$.

For arbitrary multiplier choose an actual arithmetic element
$\gamma\in G_{\mathbb Q}\subset\widehat{GT}$ with
$\lambda(\gamma)=\lambda(F)$, using cyclotomic surjectivity.
The exact arithmetic square already proved in (B.C6) gives
$D_L(\gamma)=\gamma$.  Apply the multiplier-one result to
$F\gamma^{-1}$ and use the actual homomorphism property of
$D_L$ to obtain $D_L(F)=F$.  Equation (B.C3) then gives the
second equality of (B.K10).  Iteration alternates $F$ and
$cFc$, both in full GT.  The main theorem $\Delta|_{\widehat{GT}}=0$
therefore proves the final character assertion.
\end{proof}
